\section{Implementação}

\graphicspath{{img/}{../img/}{../../img/}}

\pgfplotstableread[col sep=comma]{../../results/optimized_largeP_metrics.csv}\streammetrics

\pgfplotstablegetelem{0}{rmse}\of\streammetrics
\edef\streamFullRmse{\pgfplotsretval}
\pgfplotstablegetelem{1}{rmse}\of\streammetrics
\edef\streamIncrementalRmse{\pgfplotsretval}
\pgfplotstablegetelem{0}{coef_diff_fro}\of\streammetrics
\edef\streamFullCoefDiff{\pgfplotsretval}
\pgfplotstablegetelem{1}{coef_diff_fro}\of\streammetrics
\edef\streamIncrementalCoefDiff{\pgfplotsretval}

Os protótipos desenvolvidos ao longo da pesquisa foram consolidados no pacote
\texttt{regkit}, estruturado para atender a duas demandas complementares. A
primeira trata da reprodução fiel da formulação simbólica discutida nas seções
anteriores. Para isso, o módulo \texttt{naive} emprega a biblioteca SymPy para
gerar automaticamente todos os monômios até um grau $r$ e combina cada termo
com os versores $\hat e_d$ associados às saídas. Em seguida, as funções
lambdificadas são avaliadas ponto a ponto para montar a matriz de Gram $G$ e o
vetor $b$ do sistema normal $G\vec{a}=\vec{b}$, preservando a estrutura em
blocos descrita anteriormente.

\begin{equation*}
    G_{ij} =
    \sum_{p=1}^{N} \varphi_i(\vec{x}_p) \cdot \varphi_j(\vec{x}_p),
    \qquad
    b_i = \sum_{p=1}^{N} f(\vec{x}_p) \cdot \varphi_i(\vec{x}_p).
\end{equation*}

Embora computacionalmente onerosa, essa rota permitiu validar cada etapa da
derivação original e gerar códigos numéricos equivalentes via SymPy, que
serviram como referência para os demais experimentos.

A segunda demanda é a execução eficiente em cenários com muitos pontos
amostrais. O módulo \texttt{optimized} codifica diretamente a montagem por
blocos, utilizando a transformação \textit{PolynomialFeatures} do
\textit{scikit-learn} para obter a matriz de projeto $V \in \mathbb{R}^{N \times
Q}$. Dessa forma, o sistema normal passa a ser resolvido de maneira vetorizada,
com os $m$ problemas acoplados em uma única resolução matricial. Quando o
condicionamento de $V^{\top}V$ se deteriora, a implementação detecta o
problema e realiza uma regressão por mínimos quadrados diretamente em $V$,
garantindo estabilidade numérica\cite{burden_faires_2011}.

Para bases muito extensas, o submódulo \texttt{regression\_largeP} acumula $G$ e
$b$ iterativamente em lotes de tamanho controlado, adicionando regularização de
cresta quando necessário. Esse fluxo \textit{fora da memória} é necessário para
manipular nuvens com grande número de pontos e graus polinomiais altos sem
exceder a memória disponível. As medições consolidadas na
Tabela~\ref{tab:streaming} mostram que a versão incremental produz o mesmo erro
quadrático médio (RMSE) da solução integral,
com ambos registrando
\num[scientific-notation=true,round-mode=figures,round-precision=3]{\streamFullRmse},
e divergindo apenas por diferenças numéricas de ordem
\num[scientific-notation=true,round-mode=figures,round-precision=2]{\streamIncrementalCoefDiff}
na norma de Frobenius dos coeficientes.

\begin{table}[H]
    \centering
    \caption{Comparativo entre as variantes integral e incremental do módulo
    \texttt{regression\_largeP}.}
    \label{tab:streaming}
    \begin{tabular}{lcc}
        \hline\\
        Método & RMSE & $\|C_{\text{full}} - C_{\text{stream}}\|_F$\\\\
        \hline\\
        Integral & \num[scientific-notation=true,round-mode=figures,round-precision=3]{\streamFullRmse} & \num{\streamFullCoefDiff}\\\\
        Incremental & \num[scientific-notation=true,round-mode=figures,round-precision=3]{\streamIncrementalRmse} & \num[scientific-notation=true,round-mode=figures,round-precision=2]{\streamIncrementalCoefDiff}\\\\
        \hline
    \end{tabular}
\end{table}

Todo o código foi empacotado de forma reprodutível: os scripts
\texttt{run\_regkit\_usage.py} e \texttt{run\_benchmark\_regression.py}
geram automaticamente os conjuntos de dados, executam as regressões com
diferentes hiperparâmetros e persistem as figuras analisadas na
Seção~\ref{sec:resultados}.
